% PTE Core 模考 2B —— 听力 错题与详解·整理卷  版本 001（自包含）
% 编译: xelatex 本文件.tex   （需 Noto CJK + TeX Gyre Heros 字体）
% 本版本无外部图片，单文件即可复现。
\documentclass[11pt]{article}

% --------- 样式（原 ptestyle.sty，已内联）---------
% ============================================================
%  ptestyle.sty  —  PTE 模考错题整理卷 共享样式
%  编译引擎: XeLaTeX  (中英混排, 需要 Noto CJK 字体)
% ============================================================

\usepackage[a4paper,margin=2cm,headsep=4mm]{geometry}
\usepackage{fontspec}
\usepackage{xeCJK}
\usepackage[table]{xcolor}
\usepackage[normalem]{ulem}        % \sout 删除线
\usepackage{tabularx}
\usepackage{array}
\usepackage{ragged2e}
\usepackage{enumitem}
\usepackage[most]{tcolorbox}
\usepackage{fancyhdr}
\usepackage{graphicx}
\usepackage{pifont}                % \ding{51}=✓ \ding{55}=✗（阅读对错标注）
\usepackage{needspace}
\usepackage{tikz}

% ---------- 字体 ----------
\setmainfont{TeX Gyre Heros}            % 拉丁: Helvetica 风格(纤细清晰)
\setCJKmainfont{Noto Sans CJK SC}
\setCJKsansfont{Noto Sans CJK SC}
\newfontfamily\cjkserif{Noto Serif CJK SC}
\linespread{1.18}
\setlength{\parindent}{0pt}
\setlength{\parskip}{4pt}

% ---------- 配色 ----------
\definecolor{cWrong}{HTML}{D32F2F}     % 红  = 修改 / 错误
\definecolor{cFix}{HTML}{1B7F3B}       % 绿  = 正确建议
\definecolor{cDelete}{HTML}{1565C0}    % 蓝  = 删除
\definecolor{cInsert}{HTML}{7B1FA2}    % 紫  = 插入
\definecolor{cPartial}{HTML}{E07B00}   % 橙  = 部分得分
\definecolor{cNote}{HTML}{0277BD}      % 蓝  = 注解
\definecolor{cAccent}{HTML}{16A085}    % 主题青绿(平台主色)
\definecolor{cWrongBg}{HTML}{FBE0E0}
\definecolor{cDelBg}{HTML}{DCE8FB}
\definecolor{cInsBg}{HTML}{EEE0F5}
\definecolor{cHead}{HTML}{20303A}
\definecolor{cGrayBox}{HTML}{F4F6F7}
\definecolor{cModelBg}{HTML}{EAF7EF}
\definecolor{cYouBg}{HTML}{FFF7F5}
\definecolor{cNoteBg}{HTML}{EAF3FA}
\definecolor{cRule}{HTML}{D7DEE2}

% ---------- 行内批改宏 (复刻平台标注) ----------
% \fix{原词}{正确}  红色删除线 + 绿色上标改正  (修改/拼写/空格)
\newcommand{\fix}[2]{%
  \colorbox{cWrongBg}{\textcolor{cWrong}{\sout{#1}}}%
  \,\textsuperscript{\textcolor{cFix}{\footnotesize #2}}}
% \del{原词}  蓝色删除线  (应删除)
\newcommand{\del}[1]{%
  \colorbox{cDelBg}{\textcolor{cDelete}{\sout{#1}}}%
  \,\textsuperscript{\textcolor{cDelete}{\footnotesize 删}}}
% \ins{词}  紫色下划线  (应插入)
\newcommand{\ins}[1]{%
  \textsuperscript{\textcolor{cInsert}{\footnotesize $\wedge$}}%
  \colorbox{cInsBg}{\textcolor{cInsert}{#1}}}
% \miss{词}  灰色删除线  (口语漏读 / 未作答)
\newcommand{\miss}[1]{\textcolor{gray}{\sout{#1}}}
% \add{原词}{改正}  橙色波浪线 + "补"标签  (平台漏标、本卷补标的错误)
\definecolor{cAdd}{HTML}{E65100}
\newcommand{\add}[2]{%
  \textcolor{cAdd}{\uwave{#1}}%
  \,\textsuperscript{\textcolor{cAdd}{\footnotesize 补：#2}}}
% 高亮(口语发音): 好/中/差
\newcommand{\spk}[2]{\textcolor{#1}{#2}}

% ---------- 评分单元格着色 ----------
\newcommand{\sfull}[1]{\textcolor{cFix}{\bfseries #1}}     % 满分
\newcommand{\spart}[1]{\textcolor{cPartial}{\bfseries #1}} % 部分
\newcommand{\szero}[1]{\textcolor{cWrong}{\bfseries #1}}   % 零/低分

% ---------- 分数小标签 ----------
\newtcbox{\chip}[1][cAccent]{on line, boxsep=2pt, left=4pt, right=4pt, top=1pt,
  bottom=1pt, arc=3pt, colback=#1, colframe=#1, coltext=white, nobeforeafter,
  fontupper=\bfseries\footnotesize}

% ---------- 题块小标题 ----------
\newcommand{\ptesub}[1]{%
  \par\needspace{2\baselineskip}\smallskip
  {\color{cAccent}\rule[-0.2ex]{3pt}{1.05em}}\hspace{6pt}%
  {\bfseries\color{cHead}#1}\par\nopagebreak\smallskip}

% ---------- 题块外框 ----------
% \begin{qblock}{标号·题型}{右上信息(得分/用时)} ... \end{qblock}
\newtcolorbox{qblock}[2]{
  breakable, enhanced, arc=2.5pt,
  colback=white, colframe=cRule, boxrule=0.6pt,
  left=11pt, right=11pt, top=8pt, bottom=10pt,
  toptitle=3pt, bottomtitle=3pt, lefttitle=11pt, righttitle=11pt,
  colbacktitle=cHead, coltitle=white, fonttitle=\bfseries,
  title={#1\hfill\normalfont\bfseries\small #2}}

% ---------- 文本块: 题干 / 范文 / 你的作答 / 建议 ----------
\newtcolorbox{promptbox}{enhanced, breakable, colback=cGrayBox, colframe=cGrayBox,
  arc=2pt, left=8pt, right=8pt, top=6pt, bottom=6pt, boxrule=0pt}
\newtcolorbox{modelbox}{enhanced, breakable, colback=cModelBg, colframe=cModelBg,
  borderline west={3pt}{0pt}{cFix}, arc=1pt, left=8pt, right=8pt, top=6pt, bottom=6pt, boxrule=0pt}
\newtcolorbox{youbox}{enhanced, breakable, colback=cYouBg, colframe=cYouBg,
  borderline west={3pt}{0pt}{cWrong}, arc=1pt, left=8pt, right=8pt, top=6pt, bottom=6pt, boxrule=0pt}
\newtcolorbox{notebox}{enhanced, breakable, colback=cNoteBg, colframe=cNoteBg,
  borderline west={3pt}{0pt}{cNote}, arc=1pt, left=8pt, right=8pt, top=6pt, bottom=6pt, boxrule=0pt}

% ---------- 评分表 ----------
% 用法见正文; 这里给表头宏
\newcommand{\scorehead}{%
  \rowcolors{2}{cGrayBox}{white}
  \arrayrulecolor{cRule}}

% ---------- 章节大标题 ----------
\newcommand{\ptesection}[3]{% #1=中文名 #2=英文名 #3=该项分数
  \addcontentsline{toc}{section}{#1}%
  \needspace{6\baselineskip}
  \begin{tcolorbox}[enhanced, colback=cAccent, colframe=cAccent, sharp corners,
     left=14pt, right=14pt, top=10pt, bottom=10pt, boxrule=0pt]
    {\fontsize{20}{24}\selectfont\bfseries\color{white} #1\ \ }%
    {\large\color{white!85} #2}\hfill
    {\large\color{white}本项得分\ \ }{\fontsize{22}{24}\selectfont\bfseries\color{white} #3}%
    \ {\color{white!85}/ 90}
  \end{tcolorbox}}

% ---------- 题型小节分隔 (口语 RA/RS/DI/RTS/ASQ 等) ----------
\newcommand{\ptetask}[2]{%
  \addcontentsline{toc}{subsection}{#1\ #2}%
  \needspace{4\baselineskip}\medskip
  \begin{tcolorbox}[enhanced, colback=cHead, colframe=cHead, sharp corners,
     left=12pt, right=12pt, top=5pt, bottom=5pt, boxrule=0pt]
    {\large\bfseries\color{white} #1}\ \ {\color{white!80} #2}
  \end{tcolorbox}\smallskip}

% ---------- 口语 AI 语音识别着色 (复刻平台: 优/良/差 + 停顿) ----------
\definecolor{cSpkGood}{HTML}{2E9E5B}   % 优 绿
\definecolor{cSpkOk}{HTML}{C77F00}     % 良 黄/琥珀
\definecolor{cSpkBad}{HTML}{E0352B}    % 差 红
\newcommand{\sg}[1]{\textcolor{cSpkGood}{#1}}   % 优
\newcommand{\sy}[1]{\textcolor{cSpkOk}{#1}}     % 良
\newcommand{\sr}[1]{\textcolor{cSpkBad}{#1}}    % 差
\newcommand{\smiss}[1]{\textcolor{cSpkBad}{\sout{#1}}}  % 漏读 (红删除线)
\newcommand{\pse}{{\color{gray}\,/\,}}          % 停顿
\newcommand{\asrlegend}{{\footnotesize\color{gray}AI 语音识别\quad
  {\color{cSpkGood}\textbullet}\,优\ \ {\color{cSpkOk}\textbullet}\,良\ \
  {\color{cSpkBad}\textbullet}\,差\ \ {\color{gray}/}\,停顿\ \
  {\color{cSpkBad}\sout{词}}\,漏读}}
\newtcolorbox{asrbox}{enhanced, breakable, colback=white, colframe=cRule,
  boxrule=0.6pt, arc=1pt, left=8pt, right=8pt, top=6pt, bottom=6pt}

% ---------- 中文翻译行 ----------
\newcommand{\zh}[1]{{\footnotesize\textcolor{cNote}{\textbf{译}}\ \textcolor{gray}{#1}}}

% ---------- 阅读填空 / 选择 对错标注 ----------
\newcommand{\rok}[1]{{\color{cFix}\ding{51}\,#1}}                 % 选对(绿勾 + 词)
\newcommand{\rno}[2]{{\color{cWrong}\ding{55}\,\sout{#1}}\,{\footnotesize\color{cFix}(答案:\,#2)}}  % 选错: 你的(红删) + (答案:正确)
\newcommand{\rblank}[1]{{\color{cFix}\underline{\,#1\,}}}         % 直接给空的正确答案(无你的作答时)

% ---------- 听力专用 ----------
\newcommand{\hiw}[2]{{\color{cWrong}\uline{#1}}\,{\footnotesize\color{cFix}(听:#2)}} % HIW: 屏幕错词(红下划线)+ 实际读音
\newcommand{\lhit}[1]{{\color{cFix}#1}}                           % WFD: 命中的词(绿)
\newcommand{\lmiss}[1]{{\color{cWrong}\sout{#1}}}                 % WFD: 多余/错误的词(红删, 1B 配色)
% ---- WFD 逐词批改 · 本套 2B 平台配色（与 1B 相反）：绿=命中 / 红括号=漏写 / 灰删除线=多写或听错 ----
% 说明: 2B 平台把"漏写"标红(加括号)、"多写/听错"标灰删除线; 本卷忠于原图按此复刻。
\newcommand{\womit}[1]{{\color{cWrong}(#1)}}                      % WFD: 漏写(正确答案有、你没写)——红色括号
\newcommand{\wbad}[1]{{\color{gray}\sout{#1}}}                    % WFD: 多写/听错(你写了但错)——灰删除线
\newcommand{\wfdlegend}{{\footnotesize\color{gray}WFD 配色：\
  {\color{cFix}命中}\ \ {\color{cWrong}(漏写)}\ \ {\color{gray}\sout{多写/听错}}}}

% ---------- 颜色图例 ----------
\newcommand{\ptelegend}{%
  \begin{tcolorbox}[enhanced, colback=cGrayBox, colframe=cRule, boxrule=0.5pt,
    arc=2pt, left=10pt, right=10pt, top=6pt, bottom=6pt]
  {\bfseries\color{cHead}颜色说明\quad}
  \fix{改}{正}~修改/拼写\quad
  \del{删}~应删除\quad
  \ins{插}~应插入\quad
  \add{补}{改正}~平台漏标(本卷补标)\quad
  {\color{cFix}\rule{1em}{0.6em}}~范文/正确\quad
  {\color{cNote}\rule{1em}{0.6em}}~AI 建议
  \end{tcolorbox}}

% ---------- 页眉页脚 ----------
\pagestyle{fancy}
\fancyhf{}
\renewcommand{\headrulewidth}{0.4pt}
\renewcommand{\footrulewidth}{0pt}
\fancyhead[L]{\footnotesize\color{cHead}\bfseries PTE Core 模考 2B}
\fancyhead[R]{\footnotesize\color{gray} 错题与详解整理卷}
\fancyfoot[C]{\footnotesize\color{gray}\thepage}

\begin{document}

% --------- 封面（原 cover.tex）---------
% ---------- 封面 / 成绩总览 ----------
\definecolor{mListen}{HTML}{1F3A5F}
\definecolor{mRead}{HTML}{C2D70B}
\definecolor{mSpeak}{HTML}{4D4D4D}
\definecolor{mWrite}{HTML}{A01B7D}

\begin{titlepage}
\thispagestyle{empty}
\centering
\vspace*{1.4cm}

{\fontsize{30}{36}\selectfont\bfseries\color{cHead} PTE Core 模考 2B}\\[4mm]
{\fontsize{17}{20}\selectfont\color{cAccent}\bfseries 错题与详解 · 整理卷}\\[3mm]
{\color{gray} APEUni / 猩际 PTE 模考　·　提交时间 2026-06-28 11:49}\\[1.2cm]

% ---- 总分 ----
\begin{tcolorbox}[width=0.62\textwidth, halign=center, enhanced,
  colback=cAccent, colframe=cAccent, arc=6pt, boxrule=0pt, top=10pt, bottom=10pt]
  {\color{white}\large 总分 Overall}\\[2pt]
  {\fontsize{46}{50}\selectfont\bfseries\color{white} 50}
\end{tcolorbox}

\vspace{1.0cm}

% ---- 四项分数条形图 ----
\begin{tcolorbox}[width=0.84\textwidth, enhanced, colback=white, colframe=cRule,
  boxrule=0.8pt, arc=4pt, left=18pt, right=18pt, top=14pt, bottom=14pt]
\setlength{\tabcolsep}{6pt}\renewcommand{\arraystretch}{1.7}
\sffamily
\begin{tabular}{@{}r@{\hskip 8pt}l@{\hskip 10pt}l@{}}
 \bfseries Listening 听力 & {\color{mListen}\rule{5.22cm}{0.42cm}} & {\bfseries\large 47}\\
 \bfseries Reading 阅读   & {\color{mRead}\rule{5.11cm}{0.42cm}}   & {\bfseries\large 46}\\
 \bfseries Speaking 口语  & {\color{mSpeak}\rule{4.00cm}{0.42cm}}  & {\bfseries\large 36}\\
 \bfseries Writing 写作   & {\color{mWrite}\rule{7.89cm}{0.42cm}}  & {\bfseries\large 71}\\
\end{tabular}\\[2mm]
{\footnotesize\color{gray} 满分 90　|　刻度: 条长 = 得分 / 90}
\end{tcolorbox}

\vfill
\begin{flushleft}\small\color{gray}
\textbf{说明:}本卷由本次模考的题目与"详解"截图逐题誊写、重新排版而成,完整保留每道题的
\emph{题干 / 原文、范文、你的作答、逐项评分、彩色批改与 AI 中文建议}。\\
所有错误按平台标注用颜色标出(见下方图例)。原始"详解"PDF 不可复制、仅截图,本卷内容均为人工对照誊录,若与原图有出入以原图为准。
\end{flushleft}
\vspace{4mm}
\ptelegend
\end{titlepage}

% --------- 听力正文（原 sec_listening_body.tex）---------
% ===================== 听力 Listening =====================
\ptesection{听力}{Listening · 7 题型}{47}

\vspace{1mm}
{\small\color{gray} 本部分共 \textbf{15 题}：\textbf{SST 概述口语}Q1、\textbf{MCM-L 多选}Q2、\textbf{FIB-L 听力填空（打字）}Q3--5、\textbf{MCS-L 单选}Q6、\textbf{SMW 选词}Q7--8、\textbf{HIW 找错词}Q9--11、\textbf{WFD 听写句子}Q12--15。每题保留\textbf{录音原文}（听力口语稿，含口误照原样誊录）、你的作答与平台评分，并附\textbf{中文翻译}。
标注约定：FIB-L 空处 \rok{绿勾}=对、{\color{cWrong}\ding{55}}\,红叉=错并给答案；HIW 屏幕错词用{\color{cWrong}\uline{红下划线}}标出并括注\,{\footnotesize\color{cFix}(听:实际读音)}；WFD 逐词批改本套 2B 配色为 \wfdlegend 。
本套听力总分 \textbf{47 / 90}。}
\vspace{1mm}

{\small\color{gray}\textbf{说明：}SST 的\textbf{录音原文}此前 000 版复盘截图未提供、暂缺；\textbf{本版本（001）已由考生在平台「查看原题」取得并补入}。其余题型的录音 Transcript 复盘页均直接给出，已照抄。}
\vspace{2mm}

\newcolumntype{Z}{>{\RaggedRight\arraybackslash}X}

% ---------------- Q1 ----------------
\begin{qblock}{Q1\quad SST (C) \#5\ \ Preserving Biodiversity}{概述口语 · 评分 \textcolor{cAccent}{\bfseries 8.9 / 10}}
\ptesub{录音原文 Transcript}
\begin{promptbox}
Biodiversity, the variety of life on Earth, is essential for the stability and health of ecosystems and our own survival. The loss of biodiversity is alarming. It's not just about losing a species; it's about destabilizing entire ecosystems. For example, the decline in bee populations affects pollination, which is crucial for food crops. This shows how interconnected and dependent we are on biodiversity for food security, medicine, and ecological balance. Another vital aspect is the genetic diversity within species, key for adaptation to changing environments. For instance, diverse genetic traits in crops like rice can make them more resilient to climate change impacts. We must recognize that preserving biodiversity isn't just an environmental issue; it's a necessity for human well-being. Actions such as habitat preservation, sustainable agriculture, and combating climate change are vital steps. By protecting biodiversity, we safeguard our future, ensuring that we continue to thrive in a balanced, diverse natural world.
\end{promptbox}
\zh{生物多样性——地球上生命的多样性——对生态系统的稳定与健康、以及我们自身的生存都至关重要。生物多样性的丧失令人警醒。这不仅是失去某个物种，更是在动摇整个生态系统。例如，蜂群数量的下降会影响授粉，而授粉对粮食作物至关重要。这说明我们在粮食安全、医药与生态平衡方面，与生物多样性是何等相互关联、彼此依赖。另一个重要方面是物种内部的遗传多样性，它是适应环境变化的关键。例如，像水稻这样的作物中多样的遗传特性，能使其对气候变化的影响更具韧性。我们必须认识到，保护生物多样性不只是一个环境问题，更是人类福祉的必需。诸如栖息地保护、可持续农业与应对气候变化等行动，都是至关重要的步骤。通过保护生物多样性，我们守护自己的未来，确保能在一个平衡、多样的自然世界中持续繁荣。}
\ptesub{你的概述 Your Summary}
\begin{youbox}
In the university, we should save food and protect the \fix{environmet}{environment}. Climate change will impact human beings, and many \fix{food}{foods} like rice and so on.
\end{youbox}
\zh{在大学里，我们应当节约粮食、保护环境。气候变化将影响人类，以及许多像稻米这样的粮食等等。}
{\footnotesize\color{gray}（平台彩色标注：environmet —— 拼写错误，应为 environment；food —— 应为复数 foods。）}
\ptesub{AI 评分明细（满分 10）}
\begin{notebox}
{\renewcommand{\arraystretch}{1.35}
\begin{tabularx}{\linewidth}{@{}l c Z@{}}
\textbf{单项} & \textbf{得分} & \textbf{AI 建议}\\[2pt]
Content（内容）& \sfull{2 / 2} & 内容很棒。\\
Form（形式）& \sfull{2 / 2} & Form 很棒，这是必拿分数。\\
Grammar（语法）& \spart{1.5 / 2} & 点击有颜色标注的文字查看语法错误细节。\\
Spelling（拼写）& \sfull{2 / 2} & 虽然你拼写拿了满分，但还是犯了 1 个拼写错误，下次注意避免。\\
Vocabulary（词汇）& \spart{1.4 / 2} & PTE 对词汇不要求很难，但需要恰当。\\
\end{tabularx}}
\smallskip
\textbf{本题分值}\ 10 分\qquad\textbf{你的总得分}\ \textcolor{cAccent}{\bfseries 8.9 / 10}
\end{notebox}
\smallskip
{\footnotesize\color{gray}（答题用时 05:46；失分主要在 Grammar 与 Vocabulary。）}
\end{qblock}

% ---------------- Q2 ----------------
\begin{qblock}{Q2\quad MCM-L \#119\ \ Weight Gain}{多选 · 评分 \textcolor{cAccent}{\bfseries 2 / 2}}
\ptesub{录音原文 Transcript}
\begin{promptbox}
Both men and women gained weight but when it came to race, black women had the most rapid weight gain, followed by white women and then black and white men. The weight gain was just a few pounds---but even a slightly higher BMI is associated with weight-related health issues. Several reasons exist for the weight discrepancy between the single and married people. For example, married men and women may be less concerned about their body weight because they're no longer actively seeking a mate. Plus, married people have a regular dining partner, possibly leading to more meals. On the single side, those who are widowed or have gone through a divorce may lose weight due to stress. So while this news may not be good for the smug marrieds, it may be welcomed by Bridget Jones. The single, and very weight-conscious, Jones may actually have had the easier path to staying thin than her married friends.
\end{promptbox}
\zh{男性和女性都会增重，但就种族而言，黑人女性增重最快，其次是白人女性，再次是黑人和白人男性。增重只有几磅——但即便 BMI 略高，也与体重相关的健康问题有关。单身者与已婚者之间体重差异有几个原因。例如，已婚男女可能不太在意体重，因为他们不再积极寻找伴侣；加上已婚者有固定的用餐伙伴，可能导致进餐更多。另一方面，丧偶或离婚的单身者可能因压力而体重下降。所以这条新闻对自得的已婚者也许不是好消息，却可能受到 Bridget Jones 的欢迎——单身且非常在意体重的 Jones 也许比已婚的朋友更容易保持苗条。}
\ptesub{题目 Question}
\begin{notebox}Why are married people more likely to gain weight?\end{notebox}
\ptesub{选项 Choices（正确项绿色加粗）}
{\small
{\color{cFix}\textbf{A) Because they have regular dining partners and probably have more meals.}}\\
B) Because people's life is less stressful after marriage.\\
{\color{cFix}\textbf{C) Because they don't need to find a mate anymore and care less about their weight.}}\\
D) Because they are happier after marriage and eat more.\\
E) Because their partners require them to eat more.}
\par\smallskip
\textbf{正确答案}\ \textcolor{cFix}{\bfseries A, C}\qquad\textbf{你的答案}\ \textcolor{cFix}{\bfseries A, C}\qquad\textbf{本题得分：}\ {\color{cAccent}\bfseries 2 / 2}\quad{\footnotesize\color{gray}（用时 01:06，全对）}
\end{qblock}

% ---------------- Q3 ----------------
\begin{qblock}{Q3\quad FIB-L \#166\ \ Cultural Heritage}{听力填空（打字）· 评分 \textcolor{cAccent}{\bfseries 0 / 4}}
\ptesub{录音原文 Transcript（含填空作答）}
\begin{promptbox}
All around the world, significant parts of our cultural heritage are \rno{strenthen}{threatened} by pollution, neglect, \rno{careless}{carelessness} and greed. In learning the importance of our history, we come to understand the need to protect significant \rno{construction}{remains} from the past so that future \rno{construction}{generations} can come to understand their heritage.
\end{promptbox}
\zh{世界各地，我们文化遗产的重要部分正受到污染、忽视、疏忽大意与贪婪的威胁。在认识历史重要性的过程中，我们逐渐明白必须保护来自过去的重要遗存，好让后代也能理解他们的遗产。}
\par\smallskip
\textbf{你的作答：}\ {\color{cWrong}strenthen, careless, construction, construction}\qquad\textbf{本题得分：}\ {\color{cAccent}\bfseries 0 / 4}\quad{\footnotesize\color{gray}（用时 01:04；四空全错，正确为 threatened / carelessness / remains / generations）}
\end{qblock}

% ---------------- Q4 ----------------
\begin{qblock}{Q4\quad FIB-L \#208\ \ Well-being}{听力填空（打字）· 评分 \textcolor{cAccent}{\bfseries 1 / 3}}
\ptesub{录音原文 Transcript（含填空作答）}
\begin{promptbox}
Life in the UK 2012 provides a unique overview of well-being in the UK today. The report is the first snapshot of life in the UK to be \rno{diliver}{delivered} by the Measuring National Well-being program and will be updated and published annually. Well-being is discussed in terms of the economy, people and the environment. Information such as the \rno{unemployed}{unemployment} rate or number of crimes against the person are presented alongside data on people's thoughts and feelings, for example, \rok{satisfaction} with our jobs or leisure time and fear of crime. Together, a richer picture on `how society is doing' is provided.
\end{promptbox}
\zh{《2012 年英国生活》对当今英国的福祉作了独特的总览。该报告是首份由"国家福祉测量"项目发布的英国生活快照，并将每年更新发布。福祉从经济、民众与环境三方面来讨论。诸如失业率或针对个人的犯罪数量等信息，会与人们的想法和感受数据一并呈现，例如对工作或闲暇时间的满意度、对犯罪的担忧等。如此便提供了一幅关于"社会运行得如何"的更丰富图景。}
\par\smallskip
\textbf{你的作答：}\ {\color{cWrong}diliver}, {\color{cWrong}unemployed}, satisfaction\qquad\textbf{本题得分：}\ {\color{cAccent}\bfseries 1 / 3}\quad{\footnotesize\color{gray}（用时 01:09；仅第 3 空 satisfaction 正确，diliver→delivered、unemployed→unemployment）}
\end{qblock}

% ---------------- Q5 ----------------
\begin{qblock}{Q5\quad FIB-L \#95\ \ Oceanographer}{听力填空（打字）· 评分 \textcolor{cAccent}{\bfseries 0 / 4}}
\ptesub{录音原文 Transcript（含填空作答）}
\begin{promptbox}
For many years the favorite horror story about \rno{abrot}{abrupt} climate change was that a shift in ocean currents could \rno{directly}{radically} cool Europe's climate. These currents, called the overturning \rno{servculation}{circulation} bring warm water and warm temperatures north from the equator to Europe. Susan Loosier, an \rno{profossor}{oceanographer} at Duke University, says scientists have long worried that this ocean circulation could be disrupted.
\end{promptbox}
\zh{多年来，关于突变性气候变化最受关注的"恐怖故事"是：洋流的转变可能使欧洲气候急剧变冷。这些被称为"翻转环流"的洋流，把温暖的海水和温度从赤道向北带到欧洲。杜克大学的海洋学家 Susan Loosier 说，科学家长期担心这种海洋环流可能被打断。}
\par\smallskip
\textbf{你的作答：}\ {\color{cWrong}abrot, directly, servculation, profossor}\qquad\textbf{本题得分：}\ {\color{cAccent}\bfseries 0 / 4}\quad{\footnotesize\color{gray}（用时 01:16；四空全错，正确为 abrupt / radically / circulation / oceanographer）}
\end{qblock}

% ---------------- Q6 ----------------
\begin{qblock}{Q6\quad MCS-L \#127\ \ Autism Diagnoses}{单选 · 评分 \textcolor{cAccent}{\bfseries 0 / 1}}
\ptesub{录音原文 Transcript}
\begin{promptbox}
The first two years of a baby's life are critical for brain development. In this study, parents were given video-based tips on how best to interact with their babies according to their particular needs. By the age of three, the baby still had some developmental difficulties, but their social skills had improved, and two thirds fewer met the criteria for a diagnosis of autism spectrum disorder, compared to babies who received standard care. The research charity Autistica said the study was promising and people should not have to wait for a diagnosis for support to be provided. But the National Autistic Society said autism was not a disease that should be cured or lessened. The research team says longer term studies on more babies using the technique are now needed.
\end{promptbox}
\zh{婴儿生命的头两年对大脑发育至关重要。在这项研究中，家长得到了基于视频的指导，学习如何根据自己孩子的特定需求与之互动。到三岁时，这些婴儿仍有一些发育困难，但社交能力有所改善，符合自闭症谱系障碍诊断标准的比例比接受常规照护的婴儿少了三分之二。研究慈善机构 Autistica 表示该研究很有前景，人们不应等到确诊才获得支持。但全国自闭症协会则表示，自闭症并非一种需要治愈或减轻的疾病。研究团队称，现在需要用该方法对更多婴儿进行更长期的研究。}
\ptesub{题目 Question}
\begin{notebox}What is the attitude of the research charity Autistica with regard to their finding?\end{notebox}
\ptesub{选项 Choices（正确项绿色加粗）}
{\small
A) Negative.\\
B) Assured.\\
C) Disdainful.\\
{\color{cFix}\textbf{D) Expectant.}}}
\par\smallskip
\textbf{正确答案}\ \textcolor{cFix}{\bfseries D}\qquad\textbf{你的答案}\ {\color{cWrong}\bfseries A}\qquad\textbf{本题得分：}\ {\color{cAccent}\bfseries 0 / 1}\quad{\footnotesize\color{gray}（用时 00:59；Autistica 认为研究"有前景"，态度是期待 Expectant）}
\end{qblock}

% ---------------- Q7 ----------------
\begin{qblock}{Q7\quad SMW \#116\ \ Filtering Tools}{选词（结尾被哔替换）· 评分 \textcolor{cAccent}{\bfseries 0 / 1}}
\ptesub{录音原文 Transcript（结尾词被「哔」替换）}
\begin{promptbox}
One of the biggest surprises of the Internet is we thought when the Internet came out it would provide everyone new fundamental tools to learn about the world around them, and it would actually remove the censorship and information that we get. We basically all have our own responsibility to become our own filters. We have, and that's why it seems to me the educational system has to provide those tools those filtering tools. And that's one of the reasons why I think science is so important because of its built-in filtering tools. It recognizes that we all want to believe and as Richard Feynman said, ``The person you have to question most is yourself.'' Knowing that you want to believe when you read something that validates your beliefs, you should be skeptical of it and your beliefs and you should look out for other sources. So we have to train people that one source on the Internet is not good enough you have to search broadly to see if it's \rule[0.3ex]{1.2cm}{0.4pt}\ \textit{[beep]}.
\end{promptbox}
\zh{互联网最大的意外之一是：我们原以为互联网问世会给每个人提供了解世界的全新基础工具，并真正消除我们所获信息中的审查。基本上我们都有责任成为自己的"过滤器"。正因如此，在我看来教育体系必须提供这些过滤工具。这也是我认为科学如此重要的原因之一——它自带过滤工具：它认清我们都"想要相信"，正如理查德·费曼所说："你最需要质疑的人是你自己。"当你读到某些印证你既有信念的东西时，要意识到自己"想相信"，从而对它、对自己的信念保持怀疑，并去寻找其他来源。所以我们必须训练人们：网上单一来源并不可靠，你必须广泛检索，看看它是否经得起＿＿＿（被验证）。}
\ptesub{选项 Choices（正确项绿色加粗）}
{\small
A) skeptical\\
B) important\\
{\color{cFix}\textbf{C) validated}}}
\par\smallskip
\textbf{正确答案}\ \textcolor{cFix}{\bfseries C}\qquad\textbf{你的答案}\ {\color{cWrong}\bfseries B}\qquad\textbf{本题得分：}\ {\color{cAccent}\bfseries 0 / 1}\quad{\footnotesize\color{gray}（用时 01:06；句末应为 validated）}
\end{qblock}

% ---------------- Q8 ----------------
\begin{qblock}{Q8\quad SMW \#21\ \ Crime}{选词（结尾被哔替换）· 评分 \textcolor{cAccent}{\bfseries 1 / 1}}
\ptesub{录音原文 Transcript（结尾词组被「哔」替换）}
\begin{promptbox}
Why do we have crime? When will it all stop? IT's sad that there is so much crime in our society. It hurts so many people. Most people in the world just want to live happily and be good neighbors. Why do some people turn to crime? Money is a big reason. Many criminals pickpocket, steal, kidnap or even kill people to get money. There are many terrible crimes in the world. Perhaps the worst is ethnic cleansing. This is a crime against humanity. Many people are killed because of their color or religion. People who commit this crime really go to prison. Have you ever been a victim of crime? What do you think we need to do to reduce crime rates? Perhaps you should run \rule[0.3ex]{1.2cm}{0.4pt}\ \textit{[beep]}.
\end{promptbox}
\zh{我们为什么会有犯罪？它何时才会停止？我们的社会有这么多犯罪，令人难过，伤害了许多人。世界上大多数人只想快乐地生活、做好邻居。为什么有些人走向犯罪？金钱是一大原因。许多罪犯为了钱去扒窃、偷盗、绑架，甚至杀人。世上有许多可怕的罪行，也许最严重的是种族清洗——这是反人类罪，许多人因肤色或宗教被杀害，犯下此罪的人真的会进监狱。你曾经是犯罪的受害者吗？你认为我们需要做什么来降低犯罪率？也许你应该跑去＿＿＿（向你的政府）。}
\ptesub{选项 Choices（正确项绿色加粗）}
{\small
{\color{cFix}\textbf{A) to your government}}\\
B) to your family\\
C) to your boss\\
D) to your friends}
\par\smallskip
\textbf{正确答案}\ \textcolor{cFix}{\bfseries A}\qquad\textbf{你的答案}\ \textcolor{cFix}{\bfseries A}\qquad\textbf{本题得分：}\ {\color{cAccent}\bfseries 1 / 1}\quad{\footnotesize\color{gray}（用时 01:06，正确）}
\end{qblock}

% ---------------- Q9 ----------------
\begin{qblock}{Q9\quad HIW \#7\ \ Light}{找错词 · 评分 \textcolor{cAccent}{\bfseries 3 / 5}}
\ptesub{录音原文 Transcript（屏幕错词\,{\color{cWrong}\uline{红下划线}}\,，括注实际读音）}
\begin{promptbox}
It follows from the special theory of relativity that mass and energy are different manifestations of the same thing. Now \hiw{cleavages}{researchers} have stopped light cold --- and then brought it back to life --- or light --- in an \hiw{ribaldry}{entirely} different place. The Harvard researchers sent a light pulse into a sodium cloud cooled to about 460 degrees below zero. That's the \hiw{hypocritical}{theoretical} point where matter stops all movement. The light pulse was compressed by a factor of about 50 million at the extreme temperature, and it actually changed state --- sort of like frozen water. But it also split into two forms of matter. One stayed frozen in the sodium, while the other crawled along at just 200 meters per hour. That second form --- called a matter wave --- contains the exact information of the original light wave. And when researchers sent it into a second sodium cloud a fraction of a \hiw{freesia}{millimeter} away, the wave was converted back into light. And since fiber-optic \hiw{containers}{communications} use light, that means that someday we may be able to stop them, store them and turn them back on, just like a light switch.
\end{promptbox}
\zh{由狭义相对论可知，质量与能量是同一事物的不同表现形式。如今研究人员让光彻底"停下"——随后又让它"复活"，或者说在另一个完全不同的地方重现光。哈佛的研究人员把一束光脉冲送入冷却到约零下 460 度的钠云中，那是物质停止一切运动的理论临界点。在这一极端温度下，光脉冲被压缩了约 5000 万倍，并真的发生了相变——有点像结冰的水，还分裂成两种物质形态：一种冻结在钠中，另一种以每小时仅 200 米的速度缓慢爬行。后一种形态——称为"物质波"——精确保留了原始光波的信息。当研究人员把它送入第二团相距仅几分之一毫米的钠云时，物质波又被转换回光。既然光纤通信使用光，这意味着将来我们也许能像开关电灯一样，让光"停下、存储、再开启"。}
\par\smallskip
\textbf{答案：}\ 实际读音（Pronounced words）\ {\color{cFix}researchers, entirely, theoretical, millimeter, communications}\\
\textbf{应点错词}\ 共 \textbf{5} 个（已红下划线标出）\qquad\textbf{你点中}\ \textbf{3} 个：ribaldry, freesia, containers（均正确）\\[2pt]\textbf{本题得分：}\ {\color{cAccent}\bfseries 3 / 5}\quad{\footnotesize\color{gray}（用时 01:15；漏点 cleavages、hypocritical）}
\end{qblock}

% ---------------- Q10 ----------------
\begin{qblock}{Q10\quad HIW \#45\ \ Crop Rotation（原图标题处仅重复"HIW"）}{找错词 · 评分 \textcolor{cAccent}{\bfseries 4 / 6}}
\ptesub{录音原文 Transcript（屏幕错词\,{\color{cWrong}\uline{红下划线}}\,，括注实际读音）}
\begin{promptbox}
Crop \hiw{exploitation}{rotation} is a recognized way to keep soil and the food ecosystem healthy. Now, \hiw{licenses}{scientists} are saying that rotation could be a useful tool---in the sea. Researchers tracked a shallow-water, near-shore species used for food: sea cucumbers. They're easy to harvest and are fairly valuable. But those same attributes mean they're easy to overharvest. The \hiw{indulgence}{practice} has put some sea cucumber species at high risk of extinction, even in a relatively well-managed area, the Great Barrier Marine Park in Australia. In 2004 authorities split the Great Barrier Reef into 154 zones, where each zone was fished only once every three years. Fishers rotate through the region. In computer models, the researchers ran through dozens of \hiw{avocations}{simulations} of each zone, both before and after the divisions took place. The trials revealed that even with an identical and \hiw{snapped}{low-catch} allowance in all cases, the sea cucumbers would recover more fully under the rotation strategy than by harvesting simultaneously \hiw{couched}{throughout} the region. In fact, in trials that included rotations, the yield actually increased. The research is in the Proceedings of the National Academy of Sciences. The scientists say their results suggest that such rotation might be beneficial in these kinds of shallow marine regions around the world, particularly for species that are in high demand. Because you should not eat your seed corn. Or your seed sea cucumbers.
\end{promptbox}
\zh{作物轮作是保持土壤和食物生态系统健康的公认方法。如今科学家提出，轮作在海里也可能是有用的工具。研究人员追踪了一种近岸浅水的食用物种：海参。它们容易采捕、价值不低，但这些特点也意味着容易被过度捕捞。这种做法已使一些海参物种面临极高的灭绝风险，即便在管理相对完善的澳大利亚大堡礁海洋公园也是如此。2004 年管理部门把大堡礁划分为 154 个区，每个区每三年只捕捞一次，渔民在该区域内轮换作业。在计算机模型中，研究人员对每个区在划分前后跑了数十次模拟。试验显示：即使在所有情形下都给予相同且低捕捞量的配额，海参在轮作策略下的恢复也比在整个区域同时捕捞更充分；事实上，在包含轮作的试验中产量反而增加了。该研究发表于《美国国家科学院院刊》。科学家表示，结果说明这种轮作可能对世界各地这类浅海区域有益，尤其是对需求量高的物种。因为你不该吃掉留作种子的玉米——也别吃掉留作种源的海参。}
\par\smallskip
\textbf{答案：}\ 实际读音（Pronounced words）\ {\color{cFix}rotation, scientists, practice, simulations, low-catch, throughout}\\
\textbf{应点错词}\ 共 \textbf{6} 个（已红下划线标出）\qquad\textbf{你点中}\ \textbf{4} 个：indulgence, avocations, snapped, couched（均正确）\\[2pt]\textbf{本题得分：}\ {\color{cAccent}\bfseries 4 / 6}\quad{\footnotesize\color{gray}（用时 01:27；漏点 exploitation、licenses）}
\end{qblock}

% ---------------- Q11 ----------------
\begin{qblock}{Q11\quad HIW \#53\ \ Advertising}{找错词 · 评分 \textcolor{cAccent}{\bfseries 2 / 4}}
\ptesub{录音原文 Transcript（屏幕错词\,{\color{cWrong}\uline{红下划线}}\,，括注实际读音）}
\begin{promptbox}
Companies spend \hiw{intuitions}{billions} on advertising every year. But they're not just pushing products---they're selling their brand's ``personality,'' too. Think: Red Bull. What comes to mind? Most people say things like speed...power...hyper...extreme. Well, a pair of scientists wanted to see if the energy drink's alleged qualities would \hiw{contiguous}{influence} people's performance in a racecar video game---without sipping Red Bull. So they had 70 volunteers race cars with identical specs, but different paint jobs. Four with the logo and colors of a drink---Guinness, Tropicana, Coca-Cola or Red Bull---and one car just plain green. The drivers clocked similar times with most of the cars. But behind the wheel of the Red Bull car, they \hiw{analogously}{actually} drove more aggressively, scoring either incredibly fast race times, or their slowest---by driving recklessly and crashing. The study appears in the Journal of Consumer Psychology. Polled after the game, the players didn't realize the Red Bull image apparently \hiw{phosphoenolpyruvate}{influenced} their driving. Which suggests marketing doesn't just influence a brand's personality. It could be shaping our personalities, too, without our even knowing it.
\end{promptbox}
\zh{企业每年在广告上花费数十亿。但它们推销的不只是产品——还在推销品牌的"个性"。想想红牛，你会想到什么？多数人会说速度、力量、亢奋、极限之类。于是两位科学家想看看：在不喝红牛的情况下，这款能量饮料所宣称的特质是否会影响人们在赛车电子游戏中的表现。他们让 70 名志愿者驾驶规格相同、但涂装不同的赛车：四辆带某品牌的标志和配色——健力士、纯果乐、可口可乐或红牛——还有一辆只是纯绿色。多数车上车手用时相近，但坐进红牛车时，他们实际上开得更激进：要么跑出极快的成绩，要么因鲁莽驾驶撞车而跑出最慢成绩。该研究发表于《消费者心理学杂志》。游戏后被问及时，玩家并未意识到红牛形象明显影响了他们的驾驶。这说明营销不只塑造品牌个性，也可能在我们毫无察觉时塑造我们自己的个性。}
\par\smallskip
\textbf{答案：}\ 实际读音（Pronounced words）\ {\color{cFix}billions, influence, actually, influenced}\\
\textbf{应点错词}\ 共 \textbf{4} 个（已红下划线标出）\qquad\textbf{你点中}\ \textbf{2} 个：intuitions, analogously（均正确）\\[2pt]\textbf{本题得分：}\ {\color{cAccent}\bfseries 2 / 4}\quad{\footnotesize\color{gray}（用时 01:16；漏点 contiguous、phosphoenolpyruvate）}
\end{qblock}

% ---------------- Q12 ----------------
\begin{qblock}{Q12\quad WFD \#1973}{听写句子 · 评分 \textcolor{cAccent}{\bfseries 4 / 9}}
\ptesub{正确句子 Sentence}
\begin{promptbox}
You will acquire many skills during the academic studies.
\end{promptbox}
\zh{在学术学习期间，你将获得许多技能。}
\ptesub{你的作答（逐词批改）\quad{\normalfont\wfdlegend}}
\begin{youbox}
\lhit{you} \womit{will acquire} \wbad{are required} \lhit{many skills} \womit{during the academic} \wbad{to join in your acadamic study} \lhit{studies}
\end{youbox}
\smallskip
\textbf{本题得分：}\ {\color{cAccent}\bfseries 4 / 9}\quad{\footnotesize\color{gray}（用时 01:25；命中 you / many skills / studies，漏写 will acquire、during the academic，多写 are required、to join in your acadamic study）}
\end{qblock}

% ---------------- Q13 ----------------
\begin{qblock}{Q13\quad WFD \#1972}{听写句子 · 评分 \textcolor{cAccent}{\bfseries 7 / 10}}
\ptesub{正确句子 Sentence}
\begin{promptbox}
The temporary library will be closed in the winter break.
\end{promptbox}
\zh{临时图书馆将在寒假期间关闭。}
\ptesub{你的作答（逐词批改）\quad{\normalfont\wfdlegend}}
\begin{youbox}
\lhit{The temporary library will be closed} \womit{in} \wbad{at the} \womit{winter break} \wbad{10 minutes}
\end{youbox}
\smallskip
\textbf{本题得分：}\ {\color{cAccent}\bfseries 7 / 10}\quad{\footnotesize\color{gray}（用时 00:52；前半句全对，漏写 in、winter break，多写 at the、10 minutes）}
\end{qblock}

% ---------------- Q14 ----------------
\begin{qblock}{Q14\quad WFD \#1970}{听写句子 · 评分 \textcolor{cAccent}{\bfseries 7 / 9}}
\ptesub{正确句子 Sentence}
\begin{promptbox}
One student representative will be selected from each class.
\end{promptbox}
\zh{每个班将选出一名学生代表。}
\ptesub{你的作答（逐词批改）\quad{\normalfont\wfdlegend}}
\begin{youbox}
\lhit{one student} \womit{representative} \wbad{students presentation} \lhit{will be selected} \womit{from} \wbad{in} \lhit{each class}
\end{youbox}
\smallskip
\textbf{本题得分：}\ {\color{cAccent}\bfseries 7 / 9}\quad{\footnotesize\color{gray}（用时 00:47；漏写 representative、from，多写 students presentation、in）}
\end{qblock}

% ---------------- Q15 ----------------
\begin{qblock}{Q15\quad WFD \#1967}{听写句子 · 评分 \textcolor{cAccent}{\bfseries 8 / 10}}
\ptesub{正确句子 Sentence}
\begin{promptbox}
You must wear closed shoes when working in the lab.
\end{promptbox}
\zh{在实验室工作时，你必须穿包脚的鞋。}
\ptesub{你的作答（逐词批改）\quad{\normalfont\wfdlegend}}
\begin{youbox}
\lhit{You must wear} \womit{closed} \wbad{snow} \lhit{shoes when} \womit{working} \wbad{crossing} \lhit{in the lab} \wbad{crab}
\end{youbox}
\smallskip
\textbf{本题得分：}\ {\color{cAccent}\bfseries 8 / 10}\quad{\footnotesize\color{gray}（用时 00:59；漏写 closed、working，多写 snow、crossing、crab）}
\end{qblock}

\end{document}
